\documentclass[12pt, a4paper]{article}

\usepackage[margin=2.5cm]{geometry}
\usepackage{graphicx}
\usepackage{amsmath}
\usepackage{amssymb}
\usepackage{booktabs}
\usepackage{hyperref}
\usepackage{xcolor}
\usepackage{caption}
\usepackage{float}
\usepackage{parskip}
\usepackage{microtype}
\usepackage{natbib}
\usepackage{multirow}
\usepackage{setspace}
\usepackage{lineno}
\usepackage[T1]{fontenc}
\usepackage{lmodern}

\hypersetup{
    colorlinks=true, linkcolor=black,
    citecolor=blue!60!black, urlcolor=blue!60!black
}

\doublespacing
\linenumbers

% ── Title ─────────────────────────────────────────────────────────────────────
\title{\textbf{An End-to-End Spatial Transcriptomics Pipeline for the Breast
Cancer Tumour Microenvironment Integrating Bayesian Deconvolution and Graph
Neural Networks}}

\author{Klas Holmgren\\[0.4em]
\small\textit{Independent researcher}\\
\small\href{mailto:klas0holmgren@gmail.com}{klas0holmgren@gmail.com}}

\date{}

\begin{document}

\maketitle

% ── Structured abstract ───────────────────────────────────────────────────────
\noindent\textbf{Motivation:}
Spatial transcriptomics enables simultaneous profiling of gene expression and
tissue architecture, but integrating the complete analytical workflow --- from
quality control through deep learning --- across reproducible, well-documented
software remains challenging. No publicly available pipeline combines Bayesian
cell-type deconvolution, spatially variable gene detection, neighbourhood
co-localisation analysis, and graph-based tissue domain prediction in a single
reproducible framework applied to breast cancer.

\noindent\textbf{Results:}
We present an open-source six-stage Python pipeline for spatial transcriptomics
analysis of the breast cancer tumour microenvironment. The pipeline integrates
quality control, normalisation, Leiden clustering, cell-type deconvolution with
Cell2location, Moran's $I$ spatially variable gene detection, neighbourhood
enrichment analysis, and a two-layer graph convolutional network (GCN) that
predicts tissue domain identity from the spatial gene expression graph. Validated
on two independent 10x Genomics Visium datasets --- one FFPE and one
fresh-frozen --- the GCN achieves 75--90\% test accuracy, with Cancer Epithelial
domain F1 scores of 0.90--0.96. The pipeline is fully reproducible via a
\texttt{uv}-managed Python environment and executes end-to-end in approximately
4 hours on a consumer GPU.

\noindent\textbf{Availability and implementation:}
Source code, environment specification, and documentation are freely available at
\url{https://github.com/Klas96/spatial-breast-cancer} under the MIT licence.

\noindent\textbf{Contact:} \href{mailto:klas0holmgren@gmail.com}{klas0holmgren@gmail.com}

\newpage
\tableofcontents
\newpage

% ── 1. Introduction ───────────────────────────────────────────────────────────
\section{Introduction}

Breast cancer is the most prevalent malignancy in women worldwide, and patient
outcome is shaped not only by intrinsic tumour properties but by the composition
and spatial organisation of the surrounding tumour microenvironment
(TME) \citep{wu2021}. The TME is a heterogeneous ecosystem of cancer-associated
fibroblasts (CAFs), endothelial cells, tumour-infiltrating lymphocytes (TILs),
myeloid cells, and normal epithelial cells, each contributing distinct signals
that promote or restrain tumour progression.

Single-cell RNA sequencing (scRNA-seq) has transformed our understanding of
TME heterogeneity, but requires tissue dissociation and therefore loses spatial
context. Spatial transcriptomics, pioneered by \citet{stahl2016} and
commercialised by 10x Genomics as the Visium platform, addresses this gap by
measuring near-transcriptome-wide gene expression at spatially barcoded spots
on intact tissue sections, preserving tissue architecture at near-cellular
resolution ($55\,\mu$m spot diameter).

Despite growing adoption, end-to-end analysis of Visium data remains technically
demanding, requiring coordination across multiple specialised libraries. Key
analytical steps --- cell-type deconvolution, spatially variable gene (SVG)
detection, and spatial neighbourhood analysis --- are implemented in distinct
packages, and integration with graph-based deep learning for tissue domain
prediction adds further complexity. Several frameworks exist for individual steps
\citep{kleshchevnikov2022, palla2022}, but no publicly available, fully
integrated pipeline covers the complete analytical arc from raw SpaceRanger
output to graph neural network inference.

Here we present such a pipeline, applied to human breast cancer Visium data. The
six-stage workflow is implemented in Python, validated on two independent
datasets, and packaged for one-command reproducibility. Notably, we introduce a
two-layer graph convolutional network (GCN) \citep{kipf2017} that treats the
spatial tissue section as a graph and predicts tissue domain identity directly
from gene expression and spatial adjacency, achieving 75--90\% accuracy across
datasets. All code is open-source and available on GitHub.

% ── 2. Methods ────────────────────────────────────────────────────────────────
\section{Methods}

\subsection{Datasets}

\paragraph{Spatial transcriptomics.}
Two 10x Genomics Visium datasets were used. Dataset~1 is the FFPE Human Breast
Cancer dataset (SpaceRanger v1.3.0; 18,085 genes), available at
\url{https://www.10xgenomics.com/datasets}. Dataset~2 is the V1 Fresh-Frozen
Human Breast Cancer Block~A Section~1 (SpaceRanger v1.1.0). Both datasets are
publicly available and require no ethical approval for computational re-analysis.

\paragraph{scRNA-seq reference atlas.}
Cell-type deconvolution used the Wu et al.\ \citeyearpar{wu2021} breast cancer
single-cell atlas (100,064 cells; 9 major cell types: Cancer Epithelial, CAFs,
T-cells, Myeloid, B-cells, Endothelial, Normal Epithelial, PVL, Plasmablasts),
downloaded from CellXGene (GEO: GSE176078). Raw counts were accessed via
\texttt{adata.raw} and Ensembl gene IDs were mapped to HGNC symbols via
\texttt{adata.raw.var["feature\_name"]}.

\subsection{Software environment}

All analyses were implemented in Python~3.12.3, managed with
\texttt{uv}~0.11.14. Key dependencies: Scanpy~1.12.1 \citep{wolf2018},
Squidpy~1.8.1 \citep{palla2022}, AnnData~0.12.14, Cell2location~0.1.5
\citep{kleshchevnikov2022}, PyTorch~2.12.0+cu130, PyTorch
Geometric~2.7.0, scikit-learn~1.8.0, igraph~1.0.0, and
leidenalg~0.11.0 \citep{traag2019}. GPU-accelerated steps ran on an
NVIDIA GeForce RTX~4070 (12~GB VRAM, CUDA~13.2).

\subsection{Quality control and preprocessing}

Raw Visium data were loaded with \texttt{sc.read\_visium}. Mitochondrial gene
fraction was computed by flagging genes with the \texttt{MT-} prefix. Spots were
removed if total UMI counts $<1{,}000$, detected genes $<500$, or mitochondrial
fraction $>20\%$. Retained spots were normalised to $10{,}000$ counts per spot,
log-transformed, and the 3,000 most highly variable genes (HVGs) were identified
with the \texttt{seurat\_v3} flavour. Expression was z-scored and 50 principal
components (PCs) were computed.

\subsection{Dimensionality reduction and clustering}

A $k$-nearest neighbour (kNN) graph was built in PCA space ($k=15$, 30~PCs)
and a UMAP embedding computed for visualisation. Leiden community detection
\citep{traag2019} was applied at resolution 0.5 using the igraph backend
(\texttt{flavor="igraph", directed=False}).

\subsection{Cell-type deconvolution}

Cell-type proportions were estimated with Cell2location \citep{kleshchevnikov2022},
a negative binomial model that deconvolves spot expression into cell-type
contributions. Raw integer Visium counts were used (Cell2location requires
non-normalised data). A \texttt{RegressionModel} was first trained on the atlas
for 200 epochs to derive cell-type-specific expression signatures
$\hat{\mu}_{fg}$. The spatial \texttt{Cell2location} model was then fit to
the Visium data for 10,000 epochs ($N_{\text{cells per location}}=8$,
$\alpha_{\text{detection}}=20$), yielding posterior mean cell-type abundances
$w_{sf}$ per spot. Genes were restricted to the intersection of the Visium
variable genes and the atlas gene set before fitting.

\subsection{Spatially variable gene detection}

Spatial autocorrelation was quantified using Moran's $I$ \citep{moran1950},
computed for all genes via \texttt{sq.gr.spatial\_autocorr(mode="moran",
n\_jobs=4)} on the spatial neighbour graph (\texttt{coord\_type="grid"}).
Moran's $I \in [-1, 1]$, with values approaching 1 indicating strong
positive spatial autocorrelation. The top 12 SVGs were visualised.

\subsection{Neighbourhood enrichment analysis}

Each spot was assigned its dominant cell type (the cell type with the highest
posterior mean abundance $w_{sf}$). Cell-type co-localisation was assessed
using \texttt{sq.gr.nhood\_enrichment} \citep{palla2022}, which performs a
permutation test (1,000 permutations) and reports $z$-scores indicating
whether cell-type pairs co-localise more than expected by chance.

\subsection{Graph convolutional network for tissue domain prediction}

\paragraph{Graph construction.}
The tissue section was represented as an undirected graph $G=(V, E)$, where
each Visium spot $v_i \in V$ has feature vector $\mathbf{x}_i \in \mathbb{R}^{50}$
(PCA embedding) and edges $E$ are defined by spatial adjacency. Node labels
$y_i$ are the dominant cell-type assignments from Section~2.6.

\paragraph{Model.}
A two-layer GCN \citep{kipf2017} was implemented in PyTorch Geometric:
\begin{align}
\mathbf{H}^{(1)} &= \mathrm{ReLU}\!\left(\hat{A}\,\mathbf{X}\,\mathbf{W}^{(0)}\right), \\
\mathbf{H}^{(2)} &= \mathrm{ReLU}\!\left(\hat{A}\,\mathbf{H}^{(1)}\,\mathbf{W}^{(1)}\right), \\
\hat{\mathbf{Y}} &= \mathbf{H}^{(2)}\,\mathbf{W}^{(\mathrm{out})},
\end{align}
where $\hat{A}=\tilde{D}^{-1/2}\tilde{A}\tilde{D}^{-1/2}$ is the
symmetrically normalised adjacency with self-loops, hidden dimension is 64,
and dropout $p=0.3$ is applied after each convolution.

\paragraph{Training.}
Spots were split 68/12/20\% (train/val/test). The model was trained for 300
epochs with Adam ($\mathrm{lr}=10^{-3}$, weight decay $5\times10^{-4}$)
minimising cross-entropy loss.

% ── 3. Results ────────────────────────────────────────────────────────────────
\section{Results}

\subsection{Quality control}

After filtering, Dataset~1 retained \textbf{2,509 spots} (from $\sim$5,000
total) and Dataset~2 retained \textbf{3,795 spots}.
Figure~\ref{fig:qc} shows the QC metric distributions for Dataset~1.
The higher spot retention in Dataset~2 reflects better RNA preservation in
fresh-frozen tissue.

\begin{figure}[H]
    \centering
    \includegraphics[width=\linewidth]{../results/01_qc_distributions.png}
    \caption{QC metric distributions across Visium spots (Dataset~1, FFPE).
    From left to right: total UMI counts, number of detected genes, and
    percentage of mitochondrial reads. Dashed lines indicate the filtering
    thresholds applied.}
    \label{fig:qc}
\end{figure}

\subsection{Leiden clustering reveals spatially coherent transcriptional states}

Leiden clustering at resolution 0.5 produced spatially contiguous clusters
that closely align with tissue morphology visible in the H\&E image
(Figure~\ref{fig:clusters}). The UMAP embedding shows well-separated
transcriptional states, confirming that the kNN graph captures biologically
meaningful structure.

\begin{figure}[H]
    \centering
    \includegraphics[width=\linewidth]{../results/02_clusters.png}
    \caption{Left: UMAP embedding coloured by Leiden cluster. Right: spatial
    scatter overlaid on the H\&E image. Clusters are spatially coherent,
    reflecting the organised architecture of breast cancer tissue.}
    \label{fig:clusters}
\end{figure}

\subsection{Cell2location deconvolution maps nine cell types across the tissue}

Cell2location estimated the abundance of nine cell types at each Visium spot
using the Wu et al.\ \citeyearpar{wu2021} atlas as a reference.
The spatial abundance maps (Figure~\ref{fig:deconv}) reveal the expected
compartmentalisation of tumour, stromal, and immune populations: Cancer
Epithelial cells are concentrated in morphologically tumour-dense regions,
CAFs form stromal boundaries, and T-cell and Myeloid signals are enriched
at the tumour--stroma interface.

\begin{figure}[H]
    \centering
    \includegraphics[width=\linewidth]{../results/03_cell_type_maps.png}
    \caption{Spatial maps of estimated cell-type abundances ($w_{sf}$) from
    Cell2location. Warmer colours indicate higher abundance. Nine cell types
    from the Wu et al.\ 2021 atlas are shown.}
    \label{fig:deconv}
\end{figure}

\subsection{IFI6 and TTLL12 are the top spatially variable genes}

Moran's $I$ identified \textit{IFI6} as the top SVG in Dataset~1
($I=0.769$) and \textit{TTLL12} in Dataset~2 ($I=0.875$). Both genes
show strong spatial clustering localised to tumour-dense regions.
Figure~\ref{fig:svg} shows the top 12 SVGs for Dataset~1; genes with high
Moran's $I$ mark specific compartments including the tumour core, stroma,
and immune infiltrates.

\begin{figure}[H]
    \centering
    \includegraphics[width=\linewidth]{../results/04_top_svgs.png}
    \caption{Spatial expression maps of the top 12 spatially variable genes
    in Dataset~1, ranked by Moran's $I$.}
    \label{fig:svg}
\end{figure}

\subsection{Neighbourhood enrichment reveals tumour--stroma co-localisation}

Permutation-based neighbourhood enrichment analysis identified significant
co-localisation ($z>2$) between Cancer Epithelial and CAF-dominant spots
(Figure~\ref{fig:nhood}), consistent with the known role of CAFs in
supporting tumour invasion. T-cells and Myeloid cells show mutual
co-localisation, reflecting immune infiltrate clustering at the tumour margin.
Cancer Epithelial and T-cell spots are mutually exclusive ($z<-2$), suggestive
of immune exclusion — a hallmark of immune evasion in breast cancer.

\begin{figure}[H]
    \centering
    \includegraphics[width=\linewidth]{../results/05_nhood_enrichment.png}
    \caption{Left: neighbourhood enrichment $z$-score heatmap. Red indicates
    co-localisation above chance; blue indicates mutual exclusion. Right:
    spatial scatter coloured by dominant cell type.}
    \label{fig:nhood}
\end{figure}

\subsection{GCN predicts tissue domains with 75--90\% accuracy}

The GCN achieved test accuracy of \textbf{75\%} on Dataset~1 and
\textbf{90\%} on Dataset~2 (Table~\ref{tab:gnn}). The Cancer Epithelial
domain — the largest class — was predicted with F1 scores of 0.90 and 0.96
respectively, confirming that transcriptional and spatial signals are
sufficient to delineate the dominant tissue compartment. Performance on rare
classes (Endothelial, Normal Epithelial, PVL) was limited by severe class
imbalance rather than model capacity. The higher accuracy on Dataset~2
reflects both the larger spot count (3,795 vs 2,509) and the cleaner RNA
from fresh-frozen tissue.

\begin{table}[H]
    \centering
    \caption{GCN test-set classification report across both datasets.
    Support is the number of test spots per class.}
    \label{tab:gnn}
    \begin{tabular}{lcccccc}
        \toprule
        & \multicolumn{3}{c}{\textbf{Dataset 1 (FFPE)}}
        & \multicolumn{3}{c}{\textbf{Dataset 2 (Fresh-Frozen)}} \\
        \cmidrule(lr){2-4}\cmidrule(lr){5-7}
        \textbf{Cell type} & \textbf{F1} & \textbf{Supp.}
                           & \textbf{Acc.}
                           & \textbf{F1} & \textbf{Supp.}
                           & \textbf{Acc.} \\
        \midrule
        Cancer Epithelial  & 0.90 & 190 & \multirow{9}{*}{0.75}
                           & 0.96 & 655 & \multirow{9}{*}{0.90} \\
        CAFs               & 0.69 &  91 & & 0.47 &  21 & \\
        T-cells            & 0.67 &  90 & & 0.11 &  15 & \\
        Myeloid            & 0.77 &  72 & & 0.53 &  11 & \\
        B-cells            & 0.51 &  25 & & 0.52 &  26 & \\
        Plasmablasts       & 0.46 &  15 & & 0.00 &   7 & \\
        PVL                & 0.31 &  10 & & 0.00 &   8 & \\
        Endothelial        & 0.00 &   6 & & 0.50 &  11 & \\
        Normal Epithelial  & 0.00 &   3 & & 0.00 &   5 & \\
        \midrule
        Weighted avg       & 0.74 & 502 &  & 0.88 & 759 & \\
        \bottomrule
    \end{tabular}
\end{table}

\begin{figure}[H]
    \centering
    \includegraphics[width=\linewidth]{../results/06_gnn_spatial_pred.png}
    \caption{Side-by-side comparison of ground-truth dominant cell type (left)
    and GCN prediction (right) for Dataset~1. Agreement in the tumour core
    and stromal boundary demonstrates that the model learns tissue domain
    structure from gene expression and spatial adjacency.}
    \label{fig:gnn}
\end{figure}

% ── 4. Discussion ─────────────────────────────────────────────────────────────
\section{Discussion}

We have presented an integrated, reproducible pipeline for spatial
transcriptomics analysis of the breast cancer TME. The pipeline connects
established bioinformatics tools --- Scanpy, Squidpy, and Cell2location --- with
a graph neural network component, enabling an analytical arc from raw Visium
data to spatially-aware tissue domain prediction.

The cross-dataset validation is a key strength. Consistent biological signals
emerge across both sample preparation methods: Cancer Epithelial cells dominate
tumour-dense regions, CAFs form stromal boundaries, and immune cell co-localisation
at the tumour margin is reproducibly detected. The top SVGs differ between
datasets (\textit{IFI6} vs \textit{TTLL12}), reflecting genuine
biological heterogeneity between tumours rather than technical artefact, and
underscoring the value of applying the pipeline to multiple samples.

The GCN component demonstrates that spatial gene expression graphs encode
sufficient information to predict tissue domain identity without manual
annotation. The accuracy difference between datasets (75\% vs 90\%) is driven
primarily by class composition: Dataset~2 is heavily dominated by Cancer
Epithelial spots (86\% of test set), making the classification task easier.
Weighted F1 (0.74 vs 0.88) is a more informative metric under such imbalance.
Future work could address this with cost-sensitive learning or unsupervised
graph clustering \citep{kipf2017} to discover domains without predefined labels.

Several limitations merit acknowledgement. First, dominant cell-type assignment
is a simplification — spots are mixtures, and the GCN target labels therefore
carry inherent noise. Second, the reference atlas (Wu et al. 2021) covers a
specific set of breast cancer subtypes; performance on rare subtypes or
male breast cancer samples may differ. Third, the pipeline is optimised for
standard Visium ($\sim$55\,$\mu$m spots); adaptation to Visium HD (8\,$\mu$m)
or other platforms (Slide-seq, MERFISH) would require adjustment of the spatial
neighbour graph construction.

% ── 5. Conclusion ─────────────────────────────────────────────────────────────
\section{Conclusion}

We have released an open-source, end-to-end spatial transcriptomics pipeline
that integrates Bayesian cell-type deconvolution, spatial statistics, and graph
neural networks for characterising the breast cancer tumour microenvironment.
The pipeline is validated across FFPE and fresh-frozen Visium datasets,
reproducible via a single shell command, and designed for extension to
additional datasets, cancer types, and spatial platforms. The GCN component
provides a template for spatially-aware deep learning in cancer genomics,
achieving competitive accuracy without any manual spatial domain annotations.

% ── Data availability ──────────────────────────────────────────────────────────
\section*{Data availability}

All Visium datasets are publicly available from 10x Genomics. The scRNA-seq
reference atlas is available from CellXGene and GEO (accession GSE176078).
No new data were generated in this study.

% ── Code availability ──────────────────────────────────────────────────────────
\section*{Code availability}

All analysis code is available at
\url{https://github.com/Klas96/spatial-breast-cancer} under the MIT licence.
A frozen environment specification is provided via \texttt{uv.lock}.

% ── Funding ───────────────────────────────────────────────────────────────────
\section*{Funding}

This work received no specific funding.

% ── Conflict of interest ──────────────────────────────────────────────────────
\section*{Conflict of interest}

None declared.

% ── References ────────────────────────────────────────────────────────────────
\bibliographystyle{abbrvnat}
\bibliography{refs}

\end{document}
